% ============================================================
% DETAILED METHODOLOGY & SUBSYSTEM CONFIGURATIONS
% ============================================================

\section{Detailed Methodology \& Subsystem Configurations}

\subsection{System Architecture}

Describe the complete system architecture, identifying the major hardware,
software, communication, power, control, and user-interface blocks.

% ------------------------------------------------------------
% Standard Figure 1: System Block Diagram
% ------------------------------------------------------------
\begin{figure}[!htbp]
    \centering
    \IfFileExists{figures/block_diagram.pdf}{
        \includegraphics[
            width=0.92\textwidth,
            height=0.58\textheight,
            keepaspectratio
        ]{figures/block_diagram}
    }{
        \IfFileExists{figures/block_diagram.png}{
            \includegraphics[
                width=0.92\textwidth,
                height=0.58\textheight,
                keepaspectratio
            ]{figures/block_diagram}
        }{
            \IfFileExists{figures/block_diagram.jpg}{
                \includegraphics[
                    width=0.92\textwidth,
                    height=0.58\textheight,
                    keepaspectratio
                ]{figures/block_diagram}
            }{
                \fbox{\parbox[c][0.25\textheight][c]{0.85\textwidth}{
                    \centering
                    \textbf{PLACEHOLDER: SYSTEM BLOCK DIAGRAM}\\[0.5em]
                    Add \texttt{figures/block\_diagram.pdf},\\
                    \texttt{figures/block\_diagram.png}, or\\
                    \texttt{figures/block\_diagram.jpg}.
                }}
            }
        }
    }
    \caption{System Block Diagram}
    \label{fig:block_diagram}
\end{figure}

The block diagram should show the major functional blocks and the relationships
between them. Keep the diagram at system level rather than showing individual
component connections.

\subsection{Hardware / Circuit Design}

Describe the electrical architecture, power distribution, controllers,
sensors, actuators, interfaces, protection circuitry, and communication buses.

% ------------------------------------------------------------
% Standard Figure 2: Circuit Diagram
% ------------------------------------------------------------
\begin{figure}[!htbp]
    \centering
    \IfFileExists{figures/circuit_diagram.pdf}{
        \includegraphics[
            width=0.95\textwidth,
            height=0.62\textheight,
            keepaspectratio
        ]{figures/circuit_diagram}
    }{
        \IfFileExists{figures/circuit_diagram.png}{
            \includegraphics[
                width=0.95\textwidth,
                height=0.62\textheight,
                keepaspectratio
            ]{figures/circuit_diagram}
        }{
            \IfFileExists{figures/circuit_diagram.jpg}{
                \includegraphics[
                    width=0.95\textwidth,
                    height=0.62\textheight,
                    keepaspectratio
                ]{figures/circuit_diagram}
            }{
                \fbox{\parbox[c][0.28\textheight][c]{0.88\textwidth}{
                    \centering
                    \textbf{PLACEHOLDER: SYSTEM CIRCUIT DIAGRAM}\\[0.5em]
                    Add \texttt{figures/circuit\_diagram.pdf},\\
                    \texttt{figures/circuit\_diagram.png}, or\\
                    \texttt{figures/circuit\_diagram.jpg}.
                }}
            }
        }
    }
    \caption{System Circuit Diagram}
    \label{fig:circuit_diagram}
\end{figure}

The circuit diagram should provide the electrical-level representation of the
prototype. Ensure component designators, values, connectors, supply rails,
ground connections, and major signal paths are readable.

\subsection{Subsystem Configuration}

Describe each major subsystem separately.

\subsubsection{Subsystem 1: [Name]}

Describe hardware, function, inputs, outputs, interfaces, and operating logic.

\subsubsection{Subsystem 2: [Name]}

Describe hardware, function, inputs, outputs, interfaces, and operating logic.

\subsubsection{Subsystem 3: [Name]}

Describe hardware, function, inputs, outputs, interfaces, and operating logic.

\subsubsection{User Interface / Control Module}

Describe the operator interface, display, controls, commands, and feedback.

\subsection{Software / Control Logic}

Describe the software architecture, control strategy, data processing,
decision-making, communication logic, and fault-handling process.

% ------------------------------------------------------------
% Standard Pseudocode
% ------------------------------------------------------------
\subsection{Pseudocode / Algorithm}

Use pseudocode to describe the principal operating algorithm. Replace the
example below with project-specific logic.

\begin{algorithm}[H]
\caption{System Operating Algorithm}
\label{alg:system_operation}

\KwIn{System inputs / sensor data / user commands}
\KwOut{Control outputs / system status}

Initialize system\;
Initialize sensors, actuators, and communication interfaces\;

\While{system is active}{
    Read input parameters\;
    Validate input data\;
    Process input data\;

    \eIf{operating condition is normal}{
        Generate control output\;
        Update system status\;
    }{
        Execute fault-handling procedure\;
        Generate warning / alert\;
    }

    Update user interface / monitoring data\;
}

Shutdown system safely\;

\end{algorithm}

\subsection{Operating Sequence}

Describe the normal operating sequence from system initialization through
shutdown. Refer to Figure~\ref{fig:block_diagram} and
Algorithm~\ref{alg:system_operation} where appropriate.
